% !TEX root = ../tesi.tex

\section{Conclusions}

Manual firewall configuration does not scale with the size and complexity of modern networks and offers no guarantee of correctness, which motivated VEREFOO as an automated solution for configuring firewalls in complex topologies. But because modern networks are dynamic and change frequently, VEREFOO's monolithic strategy scales poorly; this motivates React-VEREFOO's incremental approach of reconfiguring only the portion of the network affected by the NSRs that actually changed.

This thesis pushed that same incremental principle to its extreme: rather than presenting React-VEREFOO with a batch of new requirements, the proposed Iterative React-VEREFOO introduces exactly one NSR at a time, always keeping each individual subproblem as small as it can possibly be, at the cost of turning one solver call into a chain of many smaller ones.

The approach was validated qualitatively on the CESNET academic topology. The resulting configuration was also compared against a single monolithic run of vanilla VEREFOO over the same requirements: the two solutions allocated the same number of firewalls and the same number of exception rules, differing only in where one firewall was placed. This shows the iterative process reaches a configuration correct and no less compact than the monolithic one, even though it never has visibility beyond the single requirement being processed at each step — though, as discussed next, this remains to be confirmed at a larger scale.

\section{Future Work}

The most important next step is a quantitative evaluation to measure the execution time across networks and NSR sets of varying scale, and evaluate if the iterative approach is faster than the monolithic baseline, complemented by a systematic comparison of final firewall and rule counts to quantify any trade-off in configuration quality.

Additionally two directions could improve the approach itself. Section~\ref{sec:tie-breaking-closer-look} showed that the order in which NSRs are introduced can affect the cost of later requirements when the solver faces a tie; a traffic-aware ordering could reduce future reconfigurations. Separately, this thesis only handles the addition of new NSRs — extending it to also support removals, and arbitrary mixes of the two, would cover the full range of changes mimic a live network needs.